\chapter{Spinors and the Dirac Field}

\section{Why a square root of Klein--Gordon?}

The Klein--Gordon equation is second order in time:
\[
  (\p_t^2-\nabla^2+m^2)\phi=0.
\]
Dirac looked for an equation first order in both time and space:
\[
  (\ii\gamma^\mu\p_\mu-m)\psi=0.
\]
For this to imply the Klein--Gordon equation, multiply by \((\ii\gamma^\nu\p_\nu+m)\):
\[
  (\ii\gamma^\nu\p_\nu+m)(\ii\gamma^\mu\p_\mu-m)\psi=0.
\]
Since \(m\) is constant, cross terms cancel, leaving
\[
  -\gamma^\nu\gamma^\mu\p_\nu\p_\mu-m^2.
\]
Because partial derivatives commute, only the symmetric part of \(\gamma^\nu\gamma^\mu\) matters.  We need
\[
  \frac12\{\gamma^\mu,\gamma^\nu\}=\eta^{\mu\nu}.
\]
Thus the gamma matrices obey the Clifford algebra
\[
  \{\gamma^\mu,\gamma^\nu\}=2\eta^{\mu\nu}.
\]
Then every component of \(\psi\) satisfies \((\Boxop+m^2)\psi=0\).

\section{Weyl spinors}

In the chiral basis,
\[
  \gamma^\mu=
  \begin{pmatrix}
  0&\sigma^\mu\\
  \bar\sigma^\mu&0
  \end{pmatrix},
\]
where
\[
  \sigma^\mu=(1,\bm\sigma),
  \qquad
  \bar\sigma^\mu=(1,-\bm\sigma).
\]
A Dirac spinor is
\[
  \psi=
  \begin{pmatrix}
  \psi_L\\
  \psi_R
  \end{pmatrix}.
\]
For \(m=0\), the Dirac equation splits:
\[
  \ii\bar\sigma^\mu\p_\mu\psi_L=0,
  \qquad
  \ii\sigma^\mu\p_\mu\psi_R=0.
\]
The two-component objects \(\psi_L,\psi_R\) are Weyl spinors.  They are the natural language of chiral gauge theories.

\section{The Dirac Lagrangian}

Define
\[
  \bar\psi=\psi^\dagger\gamma^0.
\]
The free Dirac Lagrangian is
\[
  \Lag=\bar\psi(\ii\gamma^\mu\p_\mu-m)\psi.
\]
Varying with respect to \(\bar\psi\) gives
\[
  (\ii\gamma^\mu\p_\mu-m)\psi=0.
\]
Varying with respect to \(\psi\) gives the adjoint equation.  The fields \(\psi\) and \(\bar\psi\) are treated as independent variables in the variational calculation, much as \(z\) and \(z^*\) are often treated independently in complex calculus.

\section{Plane-wave solutions and antiparticles}

Try
\[
  \psi(x)=u(p)e^{-\ii p\cdot x}.
\]
The Dirac equation becomes
\[
  (\gamma^\mu p_\mu-m)u(p)=0.
\]
There are positive-energy solutions \(u_s(p)\).  There are also solutions associated with negative frequencies, written
\[
  v_s(p)e^{+\ii p\cdot x}.
\]
The quantum field expansion is
\[
  \psi(x)=
  \sum_s\int\frac{\dd^3p}{(2\pi)^3}
  \frac1{\sqrt{2E_{\bm p}}}
  \left[
  b_s(\bm p)u_s(p)e^{-\ii p\cdot x}
  +
  d_s^\dagger(\bm p)v_s(p)e^{\ii p\cdot x}
  \right].
\]
The operator \(b^\dagger\) creates a particle; \(d^\dagger\) creates an antiparticle.  The negative-frequency modes do not mean negative-energy particles in the spectrum after quantization.

\section{Anticommutation and the Pauli principle}

Dirac fields are quantized with anticommutators:
\[
  \{b_r(\bm p),b_s^\dagger(\bm q)\}
  =
  (2\pi)^3\delta_{rs}\delta^{(3)}(\bm p-\bm q),
\]
and similarly for \(d,d^\dagger\).  Other anticommutators vanish.  Then
\[
  (b_s^\dagger(\bm p))^2=0,
\]
because
\[
  \{b_s^\dagger,b_s^\dagger\}=2(b_s^\dagger)^2=0.
\]
No two identical fermions can occupy the same state.  The deep spin-statistics theorem is beyond a first derivation, but the rough logic is this: Lorentz invariance, positive energy, and causality force integer-spin fields to commute and half-integer-spin fields to anticommute.

\section{The Dirac propagator}

The Feynman propagator for the Dirac field is
\[
  S_F(x-y)=\int\frac{\dd^4p}{(2\pi)^4}
  \frac{\ii(\slashed p+m)}{p^2-m^2+\ii\epsilon}
  e^{-\ii p\cdot(x-y)}.
\]
Here
\[
  \slashed p=\gamma^\mu p_\mu.
\]
Check the inverse:
\[
  (\slashed p-m)(\slashed p+m)
  =
  \slashed p\,\slashed p-m^2
  =
  p^2-m^2,
\]
because the antisymmetric part of \(\gamma^\mu\gamma^\nu\) drops out when multiplied by \(p_\mu p_\nu\).  Thus
\[
  (\ii\gamma^\mu\p_\mu-m)S_F(x-y)=\ii\delta^{(4)}(x-y).
\]

\section*{Exercises}
\addcontentsline{toc}{section}{Exercises}

\begin{exercise}
Using \(\{\gamma^\mu,\gamma^\nu\}=2\eta^{\mu\nu}\), prove that \((\slashed p)^2=p^2\).
\begin{answer}
\(\gamma^\mu\gamma^\nu p_\mu p_\nu=\frac12\{\gamma^\mu,\gamma^\nu\}p_\mu p_\nu=\eta^{\mu\nu}p_\mu p_\nu\).
\end{answer}
\end{exercise}

\begin{exercise}
Show that \((b^\dagger)^2=0\) follows from \(\{b^\dagger,b^\dagger\}=0\).
\begin{answer}
\(\{b^\dagger,b^\dagger\}=2(b^\dagger)^2\).
\end{answer}
\end{exercise}

